\section{Results}
\label{sec:results}

Table~\ref{tab:outcomes} summarizes primary outcomes across all nine experimental conditions.

\begin{table*}[t]
\centering
\caption{Primary outcomes by strategy condition. P2 Turns = number of Phase~2 submission-response cycles (each cycle consists of one provider submission and one insurer adjudication). Lines Req = total unique service lines requested by the provider across all submissions. Delivered = number of services ultimately rendered to the patient. Paid = number of claims reimbursed by the insurer in Phase~3 (Claims Adjudication).}
\label{tab:outcomes}
\begin{tabular}{lcccccc}
\toprule
Condition & Provider & Insurer & P2 Turns & Lines Req & Delivered & Paid \\
\midrule
CP\_CI & Cooperate & Cooperate & 2 & 2 & 2 & 2 \\
CP\_DI & Cooperate & Defect   & 4 & 4 & 0 & 0 \\
CP\_NI & Cooperate & Default  & 1 & 5 & 3 & 3 \\
DP\_CI & Defect    & Cooperate & 2 & 11 & 11 & 11 \\
DP\_DI & Defect    & Defect   & 4 & 7 & 0 & 0 \\
DP\_NI & Defect    & Default  & 3 & 13 & 9 & 9 \\
NP\_CI & Default   & Cooperate & 2 & 5 & 5 & 5 \\
NP\_DI & Default   & Defect   & 3 & 4 & 1 & 1 \\
NP\_NI & Default   & Default  & 6 & 2 & 0 & 0 \\
\bottomrule
\end{tabular}
\end{table*}

\subsection{Condition Summaries}
\label{sec:condition-summaries}

Across the four pure-strategy conditions (CP\_CI, CP\_DI, DP\_CI, DP\_DI), the cooperate and defect prompts produced the expected behavior: conditions with a cooperating insurer ended in full approval of all requested services, while conditions with a defecting insurer ended in complete denial, confirming that the simulation is responsive to strategy manipulation. The five conditions involving at least one default agent reveal a more nuanced picture.

\textbf{CP\_CI.} The provider requested 2 lines: infliximab (J1745, 3 induction infusions) and an acute viral hepatitis panel (80074). The insurer issued \texttt{pending\_info} requesting step-therapy documentation and biosimilar confirmation. The provider switched to biosimilar infliximab-axxq (Q5121), cited cytopenias (WBC 3.2$\times$10$^9$/L, platelets 123$\times$10$^9$/L) as steroid contraindication, and confirmed GI prescriber involvement. The insurer approved both lines. The Environment synthesized hepatitis panel results (1 synthesis event).

\textbf{CP\_DI.} The provider requested 4 lines: infliximab (J1745) plus individual hepatitis B serologies (HBsAg, anti-HBc, anti-HBs). The insurer denied all lines without \texttt{pending\_info}, citing unmet step-therapy. The provider resubmitted three times with the same cytopenias argument; all were rejected. By turn~3, the provider simplified to infliximab alone, which was still denied. The provider chose \texttt{ABANDON} with \texttt{NO\_TREAT}. No diagnostic tests were approved, so no Environment synthesis occurred.

\textbf{CP\_NI.} The provider requested 5 lines: infliximab-dyyb biosimilar (Q5102, 3 induction infusions), IV administration (96413), and hepatitis B serologies (87340, 86704, 86706). The insurer denied the drug and IV administration lines without \texttt{pending\_info}, citing unmet step-therapy, but approved all 3 screening labs. The provider chose \texttt{ABANDON} with \texttt{NO\_TREAT} for the denied treatment lines. Phase~2 resolved in 1 turn with no appeals. Three diagnostic lines were delivered and paid.

\textbf{DP\_CI.} The provider requested 11 lines: infliximab (J1745, 5 infusions), IV administration (96365), hepatitis B panel (3 tests), CMP (80053), CBC (85025), CRP (86140), MRI abdomen (74183), MRI pelvis (72197), and fecal calprotectin (83993). The insurer issued \texttt{pending\_info} for treatment lines while immediately approving all 9 diagnostic tests, triggering 9 Environment synthesis events (serologies, metabolic panel, CBC, CRP, MRI, fecal calprotectin). After provider documentation, the insurer approved all 11 lines, switching the drug to biosimilar infliximab-abda under code Q5105.\footnote{Q5105 corresponds to Retacrit (epoetin alfa-epbx), an erythropoiesis-stimulating agent, not an infliximab biosimilar. The insurer agent generated this incorrect code-drug mapping, an instance of LLM ``hallucination,'' where the model produces plausible but factually wrong output. The correct code for infliximab-abda (Renflexis) is Q5104; we use Q5104 pricing in the payoff analysis (Section~\ref{sec:payoff}).}

\textbf{DP\_DI.} The provider requested 7 lines: infliximab-axxq biosimilar (Q5121, 3 induction infusions), IV administration (96413), and screening labs (HBsAg, anti-HBc, anti-HBs, HCV Ab, CRP). The insurer denied all 7 lines without \texttt{pending\_info}. The provider resubmitted, then escalated to Level~1 appeal (Medical Director), then Level~2 (IRE). All denied at every level. The provider chose \texttt{ABANDON} with \texttt{NO\_TREAT}. As in CP\_DI, no Environment synthesis occurred.

\textbf{DP\_NI.} The provider requested 13 lines: infliximab-dyyb biosimilar (Q5102, 9 infusions), IV administration (96413, 96415 additional hour), diphenhydramine premedication (J1200), hepatitis B serologies (87340, 86704, 86706), CBC (85025), CMP (80053), CRP (86140), MRI abdomen (74183), MRI pelvis (72197), and fecal calprotectin (83993). The insurer denied all treatment lines (drug, IV, premedication) without \texttt{pending\_info}, citing unmet step-therapy, but approved all 9 diagnostic lines. The provider appealed through Level~1 and Level~2; treatment lines remained denied at every level. The provider chose \texttt{ABANDON} with \texttt{NO\_TREAT} for the 4 denied lines. Nine diagnostic lines were delivered and paid.

\textbf{NP\_CI.} The provider requested 5 lines: infliximab (J1745, 3 induction infusions), IV administration (96413), MRI abdomen (74183), MRI pelvis (72197), and fecal calprotectin (83993). The insurer issued \texttt{pending\_info} requesting step-therapy documentation for the drug, modified the MRI abdomen line, and approved the remaining diagnostic tests. The provider accepted the modification and supplied documentation. The insurer approved all lines. All 5 lines were delivered and paid.

\textbf{NP\_DI.} The provider requested 4 lines: infliximab-axxq biosimilar (Q5121, 3 induction infusions), IV administration (96413), MRI abdomen (74183), and MRI pelvis (72197). The insurer denied all 4 lines without \texttt{pending\_info}. The provider appealed through Level~1 (Medical Director) and Level~2 (IRE). At Level~2, the reviewer approved infliximab but upheld denial of the remaining 3 lines. The provider chose \texttt{ABANDON} with \texttt{NO\_TREAT} for the denied lines. One line was delivered and paid.

\textbf{NP\_NI.} The provider requested 2 lines: infliximab (J1745, 9 infusions) and IV administration (96413, 9 visits). The insurer denied both lines without \texttt{pending\_info}, citing unmet step-therapy. The provider resubmitted with reduced quantities, then appealed through all levels (4 appeals over 6 turns). All denied at every level. The provider chose \texttt{ABANDON} with \texttt{NO\_TREAT}. No services were delivered; this was the most protracted negotiation across all nine conditions, despite involving only two service lines, the fewest of any condition.

\subsection{Service Line Composition}
\label{sec:service-lines}

Table~\ref{tab:service-lines} breaks down requested service lines by clinical category, revealing how the defect strategy inflates request volume primarily through ancillary services rather than additional drug lines.

\begin{table*}[t]
\centering
\caption{Service lines requested by category and condition.}
\label{tab:service-lines}
\begin{tabular}{lccccccccc}
\toprule
Category & CP\_CI & CP\_DI & CP\_NI & DP\_CI & DP\_DI & DP\_NI & NP\_CI & NP\_DI & NP\_NI \\
\midrule
Treatment (drug) & 1 & 1 & 1 & 1 & 1 & 2 & 1 & 1 & 1 \\
IV administration & 0 & 0 & 1 & 1 & 1 & 2 & 1 & 1 & 1 \\
Hepatitis screening & 1 & 3 & 3 & 3 & 4 & 3 & 0 & 0 & 0 \\
Baseline labs & 0 & 0 & 0 & 3 & 1 & 3 & 0 & 0 & 0 \\
Imaging & 0 & 0 & 0 & 2 & 0 & 2 & 2 & 2 & 0 \\
Monitoring & 0 & 0 & 0 & 1 & 0 & 1 & 1 & 0 & 0 \\
\midrule
\textbf{Total} & \textbf{2} & \textbf{4} & \textbf{5} & \textbf{11} & \textbf{7} & \textbf{13} & \textbf{5} & \textbf{4} & \textbf{2} \\
\bottomrule
\end{tabular}
\end{table*}

Cooperating providers requested 2--5 service lines; defecting providers requested 7--13; default providers requested 2--5. The additional volume from defecting providers came entirely from ancillary services: IV administration, premedication, baseline laboratory panels, imaging, and monitoring tests. Cooperating and default providers focused narrowly on the biologic drug and hepatitis screening required before infliximab initiation.

\subsection{Default Agent Behavior: Cooperation versus Defection}
\label{sec:default-insurer}

The five default-strategy conditions, in which one or both agents received no explicit behavioral instructions, serve as a natural baseline that sharpens the contrast between cooperation and defection.

\paragraph{Insurer side.}
The cooperating insurer used \texttt{pending\_info} to request clarification and then approved the drug once the provider documented cytopenias as a steroid contraindication. The defecting insurer denied all requested services, both the drug and diagnostic tests, without ever requesting additional information. The default insurer fell between these poles but closer to defection: it denied the drug in all three conditions (CP\_NI, DP\_NI, NP\_NI) while approving diagnostic tests when requested (CP\_NI: 3/3 labs; DP\_NI: 9/9 diagnostics). Crucially, the default insurer never issued a \texttt{pending\_info} request, matching the defecting insurer on the very mechanism that determines whether the drug is covered. This ``selective denial,'' blocking the treatment while permitting diagnostics, directly mirrors the structure of the coverage policy: Cigna IP0660 imposes step-therapy requirements only on the biologic drug, not on diagnostic tests. Without interpretive guidance, the default insurer applied coverage policy literally, denying requests where explicit step-therapy criteria were unmet and approving those where no restrictive criteria applied. The implication is that a cooperative posture, treating ambiguity as an invitation for dialogue rather than as grounds for denial, requires explicit interpretive guidance. Without it, the LLM defaults to literal policy application, which in gray-zone cases produces denial.

\paragraph{Provider side.}
Cooperating providers requested the minimum defensible set of services (2--5 line items), while defecting providers inflated their requests with ancillary services (7--13 line items). Default providers matched cooperators in request volume (2--5 lines) but diverged in composition: all three default-provider conditions included IV administration (CPT 96413) and two included imaging (MRI abdomen and pelvis), services that cooperating providers omitted, while none included hepatitis screening or baseline labs (Table~\ref{tab:service-lines}). Measured by total reimbursement when the insurer cooperates, the default provider behaved more like a cooperator but with a slightly broader clinical interpretation. The critical difference emerged in appeal behavior: when the cooperating insurer issued \texttt{pending\_info} (NP\_CI), the default provider supplied documentation, accepted a modification, and resolved in 2 turns. When faced with denial (NP\_DI, NP\_NI), the default provider appealed through the maximum review level before abandoning. NP\_NI was the most protracted condition in the experiment (6 turns, 4 appeals), with the provider reducing infusion quantity from 9 to 3 and changing the diagnosis code to emphasize the obstructive phenotype mid-negotiation.

\paragraph{Interaction at NP\_NI.}
When neither agent received strategic guidance, the outcome collapsed to complete denial through a distinctive pathway: the default provider requested only 2 lines (drug + IV admin), both subject to step-therapy, and the default insurer denied both. The end result matches that of mutual defection (DP\_DI: zero services approved), but the mechanism differs in instructive ways: in NP\_NI, a narrow request met literal denial and triggered persistent appeals, whereas in DP\_DI, an inflated request met blanket rejection. The default--default pair thus reaches the same deadlock as mutual defection, but through bureaucratic attrition rather than overt strategic aggression. This suggests that unguided AI systems on both sides of prior authorization could reproduce the worst cooperative outcome, not through adversarial intent, but through literal policy application by the insurer and exhaustive procedural appeals by the provider.

\subsection{Utility Function and Payoff Analysis}
\label{sec:payoff}

To formalize the strategic structure in game-theoretic terms, we define utility functions (numerical measures of each agent's outcome quality) over the simulation results. Intuitively, the provider benefits from higher reimbursement and is penalized by administrative costs incurred during negotiation; the insurer benefits from denying expensive services (reducing payouts) but also incurs administrative costs. Let $R^P$ denote the total reimbursement the provider receives for approved and paid services, let $S^I$ denote the total value of services on the final submitted request (the insurer's exposure), and let $C_P = n \cdot c_P$ and $C_I = n \cdot c_I$ denote the total administrative costs for provider and insurer respectively, where $n$ is the number of Phase~2 negotiation turns and $c_P$, $c_I$ are per-turn costs. We define:
\begin{align}
    U_P &= R^P - \alpha_P \cdot C_P \label{eq:u_provider} \\
    U_I &= S^I - R^P - \alpha_I \cdot C_I \label{eq:u_insurer}
\end{align}
where $\alpha_P, \alpha_I \geq 0$ weight how much each agent values administrative burden relative to financial outcomes. The provider maximizes reimbursement net of admin burden. The insurer's utility captures the savings from denied services: $S^I - R^P$ equals the full submitted value $S^I$ when every service is denied (maximum savings) and $0$ when every service is approved (no savings). This term is reduced by administrative burden. When $\alpha = 0$, admin costs are ignored; when $\alpha = 1$, a dollar of admin cost is weighted equally to a dollar of reimbursement or savings.

We price services using CMS reimbursement rates: drug costs from CMS Average Sales Price (ASP) data for Q4~2025, laboratory costs from the CMS Clinical Lab Fee Schedule 2025, and per-turn administrative costs from the CAQH 2023 Index. The CAQH estimates are \$5.79 per transaction for providers and \$0.05 per transaction for insurers when prior authorization is conducted electronically. Four procedure codes (96365, 96413, 74183, 72197) lack direct CMS reimbursement rates and instead use benchmark estimates bounded by verified CMS Physician Fee Schedule (PFS) facility rates. These four codes account for less than 15\% of total service value in any condition, limiting the impact of estimation uncertainty on payoff calculations. Full pricing sources are in Appendix~\ref{sec:cost-data}.

We first analyze the 2$\times$2 cooperate-defect subgame. Table~\ref{tab:payoff} presents the payoff matrix at $\alpha_P = \alpha_I = 1$ (admin costs weighted at face value); we then examine how the game's structure changes under other parameterizations.

\begin{table*}[t]
\centering
\caption{Payoff matrix ($U_P$, $U_I$) at $\alpha_P = \alpha_I = 1$. Service values and admin costs in USD. Admin: $c_P = \$5.79$/turn, $c_I = \$0.05$/turn.}
\label{tab:payoff}
\begin{tabular}{lcc}
\toprule
 & \textbf{Insurer Cooperate} & \textbf{Insurer Defect} \\
\midrule
\textbf{Provider Cooperate} & (\$1,868; $-$\$0.10) & ($-$\$23; \$2,828) \\
\textbf{Provider Defect} & (\$5,157; $-$\$0.10) & ($-$\$23; \$2,274) \\
\bottomrule
\end{tabular}
\end{table*}

The insurer has a strictly dominant strategy to defect, meaning defection yields higher utility than cooperation regardless of what the provider does: $U_I^D > U_I^C$ for both provider strategies (\$2,828~$>$~$-$\$0.10 when the provider cooperates; \$2,274~$>$~$-$\$0.10 when the provider defects). Denying all services and retaining thousands of dollars in savings dominates paying out reimbursement, even after accounting for negligible admin costs. The provider has a \emph{weakly} dominant strategy to defect: strictly preferred when the insurer cooperates (\$5,157~$>$~\$1,868), tied when the insurer defects ($-$\$23~$=$~$-$\$23). The Nash equilibrium is mutual defection ($-$\$23, \$2,274).

Whether this constitutes a prisoner's dilemma depends on parameterization. A prisoner's dilemma requires that mutual cooperation leave \emph{both} players better off than mutual defection (a condition known as Pareto dominance). At $\alpha = 1$ with observed per-turn costs, the insurer \emph{prefers} mutual defection (\$2,274) to mutual cooperation ($-$\$0.10): denied-service savings dwarf the trivial admin cost. Under these parameters, the game is better characterized as asymmetric exploitation, where the insurer's dominant strategy imposes costs on the provider, rather than a classic prisoner's dilemma. Both players defect at equilibrium, but for different reasons: the insurer defects because denial savings always exceed admin costs; the provider defects because, conditional on insurer cooperation, requesting a broader service set extracts more reimbursement, and conditional on insurer defection, the outcome collapses to pure admin loss regardless of strategy. However, the prisoner's dilemma structure emerges when insurer defection becomes more costly, either by increasing $\alpha_I$ (weighting administrative burden more heavily in the utility function) or by incorporating real-world costs such as litigation risk, provider network attrition, and reputational damage among patients and employers. The sensitivity analysis below quantifies the threshold at which the game transitions from exploitation to prisoner's dilemma.

The insurer's dominant defection strategy is driven by a striking 116-fold asymmetry in per-transaction administrative costs: providers pay \$5.79 per prior authorization transaction, while insurers pay just \$0.05. The insurer pays almost nothing for additional turns, so denial is nearly free. This result formalizes a dynamic widely described in the health policy literature: prior authorization imposes administrative costs almost entirely on providers and their patients, giving insurers little financial incentive to streamline the process \citep{ama2024pa}.

\paragraph{Sensitivity to burden weight.}
The insurer's dominant defection is robust to the burden-weight parameters $\alpha_P$ and $\alpha_I$. Because insurer admin costs are so low ($c_I = \$0.05$/turn), the insurer's strategy only changes when $\alpha_I$ reaches extreme values: the insurer becomes indifferent between cooperation and defection against a defecting provider at $\alpha_I \approx S^{DD} / \Delta C_I = \$2{,}274 / (2 \times \$0.05) \approx 22{,}700$, meaning admin costs would need to be weighted ${\sim}23{,}000\times$ their dollar value to alter the insurer's incentive. Against a cooperating provider, the threshold is even higher ($\alpha_I \approx 28{,}300$). These thresholds drop dramatically if per-turn insurer costs rise to match provider costs (e.g., at $c_I = c_P = \$5.79$, the thresholds fall to $\alpha_I \approx 196$ and $\approx 244$), values at which the game transitions to a prisoner's dilemma, consistent with the coordination-failure hypothesis of \citet{moritz2025lancet}. Whether real-world insurer costs (incorporating litigation, network, and reputational factors) reach these levels is an empirical question. The interactive sensitivity analysis in the supplementary materials allows exploration of all six parameters ($R^{CC}$, $R^{DC}$, $S^{CD}$, $S^{DD}$, $\alpha_P$, $\alpha_I$).

% \subsection{LLM Stochasticity}
% \label{sec:stochastic}
%
% Large language model outputs are inherently stochastic: identical prompts and inputs can produce different outputs across runs due to sampling randomness in the model's text generation process. We observed meaningful variation in two additional runs of selected conditions (CP\_CI and CP\_DI), conducted under identical strategy settings.
%
% Under CP\_CI, the main run requested 2 lines using a bundled hepatitis panel (80074) and selected biosimilar Q5121 (Avsola), while an alternate run requested 4 lines with individual serologies (87340, 86704, 86706) and selected biosimilar Q5103 (Inflectra). Both reached full approval in 2 turns, but through different service compositions and biosimilar selections.
%
% Under CP\_DI, stochastic variation produced qualitatively different outcomes. The main run lasted 4~turns: the provider resubmitted three times, eventually simplified to infliximab alone, and chose \texttt{ABANDON} with \texttt{NO\_TREAT} (zero services delivered, provider payoff $-$\$23). An alternate run lasted only 2~turns: the provider resubmitted once with added hepatitis labs, then upon blanket denial chose \texttt{TREAT\_ANYWAY} for the labs (subsequently paid at claims) but \texttt{NO\_TREAT} for infliximab (partial delivery, provider payoff \$18). The same strategy pairing yielded payoffs of $-$\$23 versus \$18 depending solely on LLM sampling.
%
% This run-to-run variation is not merely quantitative; it affects the qualitative game-theoretic structure of the interaction. With the main run's CP\_DI payoff ($-$\$23), the provider is indifferent under insurer defection and defection is weakly dominant. With the alternate run's payoff (\$18), the provider strictly prefers cooperation under insurer defection, eliminating the provider's dominant strategy entirely and shifting the Nash equilibrium from (D,D) to (C,D). The game-theoretic classification of the interaction is itself dependent on LLM sampling randomness, a finding that underscores the need for multiple replications per condition in future work and raises questions about the reliability of single-run AI-mediated decisions in real-world prior authorization.
